\documentclass[conference]{IEEEtran}

\usepackage[T1]{fontenc}
\usepackage[utf8]{inputenc}
\usepackage{cite}
\usepackage{url}
\usepackage{amsmath,amssymb}
\usepackage{graphicx}
\usepackage{booktabs}
\usepackage{array}
\usepackage{multirow}
\usepackage{tabularx}
\usepackage[hidelinks]{hyperref}

\begin{document}

\title{Measuring and Mitigating NER Accuracy\\
Degradation in Creative Resumes}

\author{
\IEEEauthorblockN{Prayag Adhikari}
\IEEEauthorblockA{
Department of Computer Science\\
Caldwell University\\
Caldwell, NJ, USA\\
Email: padhikari1@caldwell.edu\\
Advisor: Dr. Sushil Poudel
}
}

\maketitle

\begin{abstract}
Automated resume parsers can lose or misclassify information when document extraction and named entity recognition (NER) fail. This study evaluated resume parsing across document layouts and examined whether GPT-assisted and hybrid methods improve technical-skill recovery. An initial gold-standard experiment used 20 resumes, comprising 10 standard and 10 creative layouts, annotated for PERSON, ORG, DATE, GPE, and SKILL entities, totaling 1,108 entities. Contrary to the proposed hypothesis, the baseline achieved a higher strict F1 score on creative resumes (0.329) than on standard resumes (0.223), with an overall F1 of 0.270; therefore, the small initial sample did not demonstrate systematic layout-related degradation. The study then expanded to 1,020 computer-science resumes from three public datasets. On a stratified, manually reviewed sample of 50 resumes, GPT-only skill extraction achieved strict precision, recall, and F1 scores of 0.966, 0.644, and 0.773. A raw union of rule-based V4 and GPT predictions increased recall but reduced precision, whereas a cross-validation-derived filtered hybrid achieved 0.921 precision, 0.679 recall, and the highest strict F1 of 0.782. Bootstrap analysis estimated a combined F1 improvement of 0.0089 over GPT-only, although the confidence interval slightly crossed zero. An analysis of 50 job descriptions and a four-scan Jobscan pilot supplied exploratory external comparisons. Overall, hybrid filtering improved skill recovery, but performance varied across datasets, and claims about layout bias or commercial ATS behavior require cautious interpretation.
\end{abstract}

\begin{IEEEkeywords}
applicant tracking systems, named entity recognition, resume parsing, skill extraction, large language models
\end{IEEEkeywords}

\section{Introduction}

Resumes are semi-structured documents that combine free text, headings, dates, organizations, locations, and lists of technical skills. Automated hiring platforms may extract these elements to support candidate search, filtering, and job matching. However, published guidance about preparing resumes for Applicant Tracking Systems (ATS) is inconsistent, particularly regarding formatting and keyword use \cite{lookadoo2025}. This inconsistency creates a practical problem for students and other job seekers: a resume must communicate effectively to a human reader while remaining sufficiently machine-readable for automated processing.

Natural language processing methods have been applied to resume parsing, candidate ranking, keyword matching, and job-description comparison \cite{bhatia2019resume,yadav2025}. Named Entity Recognition (NER) provides one established framework for locating and labeling spans such as people, organizations, dates, and locations \cite{conll2003}. Nevertheless, resumes differ from ordinary prose. Multi-column layouts, graphics, text boxes, embedded images, and non-linear reading order can disrupt PDF text extraction before NER is applied. Recent layout-aware research identifies document heterogeneity and reading-order reconstruction as central resume-parsing challenges \cite{zhu2025layout}. ResumeBench similarly demonstrates that structure-rich resumes remain difficult for large language models, even when outputs follow a predefined schema \cite{ling2025resumebench}. In addition to extraction reliability, automated resume screening raises fairness concerns because resumes can encode demographic information that influences downstream models \cite{parasurama2021}.

This study investigates two related but distinct technical problems. The first is whether creative or complex PDF layouts reduce accuracy in an ATS-style extraction and NER pipeline. This question is evaluated using 20 resumes, divided evenly between standard and creative or complex layouts, with 1,108 manually annotated PERSON, ORG, DATE, GPE, and SKILL entities. The second problem is whether technical-skill recovery can be improved at a larger scale by combining a precision-focused rule-based method with contextual extraction from GPT-4o-mini. This question is studied using 1,020 CS/IT resumes collected from three public datasets. Because the large corpus is stored primarily as extracted text, it supports skill-extraction and scalability analysis but not direct conclusions about visual layout.

The main contributions of this research are:

\begin{enumerate}
    \item A comparative baseline evaluation of standard and creative or complex PDF resumes using strict exact-match scoring and entity-level error analysis.
    \item A stratified 50-resume manual evaluation comparing a generalized rule-based V4 system, GPT-only extraction, a raw V4--GPT union, and a cross-validation-filtered hybrid.
    \item A full-corpus application of GPT and hybrid methods to 1,020 resumes, accompanied by dataset-specific reporting and bootstrap uncertainty estimates.
    \item An industry-keyword analysis based on 50 CS job descriptions and an exploratory four-scan comparison with Jobscan.
\end{enumerate}

The initial layout experiment did not support the proposed expectation that creative resumes would necessarily perform worse: the creative group achieved a strict F1 score of 0.329, compared with 0.223 for the standard group. This result is interpreted cautiously because the sample was small and potentially affected by differences in profession, entity density, anonymization, and vocabulary. In the final 50-resume skill evaluation, GPT-only achieved a strict F1 score of 0.773, the raw hybrid achieved 0.754, and the filtered hybrid achieved the highest strict F1 score of 0.782. The filtered hybrid showed a clear improvement for the Snehaan source, while its combined advantage over GPT-only was modest and its bootstrap confidence interval slightly crossed zero.

Accordingly, this work does not claim to reproduce any proprietary ATS or establish universal layout bias. Instead, it provides a reproducible analysis of extraction failure, NER behavior, skill-recovery tradeoffs, and dataset variation. The results emphasize that layout evaluation and large-scale skill extraction must be treated as separate experimental questions and that practical resume guidance should be based on transparent evidence rather than unsupported claims about commercial systems.

\section{Related Work}

\subsection{ATS Guidance and Automated Resume Analysis}

Applicant Tracking Systems support activities such as storing applications, searching candidate records, extracting qualifications, and comparing resumes with job requirements. However, their internal parsing and ranking mechanisms are generally not available for independent inspection. Lookadoo and Moore audited 18 business and technical communication textbooks and found substantial variation in how the texts described AI-assisted ATS, resume formatting, and keyword practices \cite{lookadoo2025}. Their findings demonstrate that common resume advice is not always consistent or sufficiently evidence-based, strengthening the need for transparent experiments that separate extraction performance from general claims about commercial hiring platforms.

Prior technical work has applied natural language processing to resume extraction and job matching. Bhatia et al.\ presented an end-to-end system that combined resume parsing with BERT-based candidate ranking against job descriptions \cite{bhatia2019resume}. Yadav et al.\ similarly described a resume-analysis system that extracted candidate information and used keyword and ATS-style matching for evaluation against job requirements \cite{yadav2025}. These systems illustrate the practical connection between accurate information extraction and downstream candidate comparison: information that is missed or mislabeled during parsing cannot be used reliably during matching.

\subsection{Named Entity Recognition and Skill Extraction}

Named Entity Recognition provides a standard framework for labeling text spans by semantic category. The CoNLL-2003 shared task helped establish entity-level evaluation for categories such as people, organizations, and locations \cite{conll2003}. Modern NLP libraries such as spaCy provide pretrained statistical pipelines that make these methods accessible for applied research \cite{spacy2020}. However, general-purpose NER categories do not directly cover the full range of resume information. Technical skills, software tools, cloud services, frameworks, and multi-word competencies are particularly domain-specific and are not represented by a standard spaCy SKILL label.

Consequently, resume parsing frequently combines statistical NER with rules, dictionaries, or contextual models. The present research follows this combined approach: spaCy identifies PERSON, ORG, DATE, and GPE entities, while technical skills are evaluated using a custom matcher, GPT-based recovery, and hybrid combinations. This design permits separate analysis of conventional named entities and the specialized skill vocabulary that is central to CS/IT resumes.

\subsection{Layout-Aware and LLM-Based Resume Parsing}

Resume layout introduces challenges that are not captured by text-only NER evaluation. Zhu et al.\ describe a layout-aware framework that combines PDF metadata, OCR content, reading-order reconstruction, and efficient language-model extraction \cite{zhu2025layout}. Their work identifies document heterogeneity, multi-column structure, and image-based content as major barriers to consistent resume parsing. This supports evaluating text extraction as a distinct stage rather than assuming that the original PDF has already been converted into a correct linear sequence.

ResumeBench provides a complementary large-scale perspective. It contains 2,500 privacy-compliant synthetic resumes across 50 templates, 30 career fields, and five languages, and evaluates 24 language models on structured resume extraction \cite{ling2025resumebench}. Its results show substantial variation among models and continued difficulty with structural alignment and semantic ambiguity. ResumeBench offers broad template and language coverage, whereas the present study contributes a smaller manually annotated experiment using real PDF layouts and a separate 1,020-resume text corpus for skill-recovery analysis.

\subsection{Fairness, Scope, and Research Gap}

Parsing performance also matters because resume representations can influence later automated decisions. Parasurama and Sedoc showed that resumes contain gender-associated information and that removing such signals creates tradeoffs between obfuscation and screening performance \cite{parasurama2021}. Although the present study does not evaluate demographic fairness or hiring decisions, this work motivates careful attention to extraction reliability, privacy, and the limits of automated evaluation.

Together, prior studies establish the importance of resume parsing, layout awareness, specialized entity extraction, and responsible interpretation. They also leave a specific methodological gap addressed here: a manually annotated comparison between standard and creative PDF layouts should not be conflated with a text-only scalability experiment. This research therefore separates its manually annotated 20-PDF layout study from its 1,020-resume technical-skill study and reports results by dataset source as well as in aggregate.

\section{Research Questions}

This study addresses four research questions spanning the 20-resume
PDF-layout experiment, the larger technical-skill evaluation, and the
exploratory external comparisons:

\begin{enumerate}
    \item \textbf{RQ1:} How does baseline ATS-style text extraction and
    named entity recognition performance differ between standard and
    creative or complex PDF resumes, both overall and across PERSON,
    ORG, DATE, GPE, and SKILL entities?

    \item \textbf{RQ2:} To what extent does GPT-assisted reconstruction
    improve strict entity recovery over the original spaCy and
    keyword-matching baseline on the 20-resume dataset, and does the
    effect differ between layout groups?

    \item \textbf{RQ3:} On the stratified, manually reviewed sample of
    50 CS/IT resumes, how do generalized V4, GPT-only, raw hybrid, and
    cross-validation-filtered hybrid skill extraction compare under
    strict and relaxed precision, recall, and F1, both overall and by
    dataset source?

    \item \textbf{RQ4:} When GPT and hybrid methods are applied to the
    full 1,020-resume corpus, what technical-skill distributions and
    industry-keyword recovery patterns emerge relative to 50 CS job
    descriptions and the exploratory Jobscan pilot?
\end{enumerate}

The original directional hypothesis associated with RQ1 predicted that
creative or complex resumes would achieve lower strict F1 scores than
standard resumes. The remaining questions are treated primarily as
comparative or exploratory analyses. RQ1 and RQ2 use the
20-PDF layout experiment, RQ3 uses manually reviewed predictions, and RQ4 is
descriptive because the full corpus does not contain complete manual
annotations or preserved visual-layout labels.

\section{Methodology}

\subsection{Original 20-Resume Layout Experiment}

The first experiment evaluated whether resume layout was associated
with differences in ATS-style text extraction and named entity
recognition. Twenty real resume PDFs were selected and manually divided
into two equally sized groups: 10 standard resumes and 10 creative or
complex resumes. Standard resumes generally followed a conventional
single-column structure, whereas the creative or complex group included
features such as multiple columns, text boxes, graphics, unusual
section placement, and non-linear reading order. These layout labels
were assigned from the original PDFs before system performance was
calculated.

Text was first extracted from each PDF using PyMuPDF. When ordinary PDF
extraction returned no usable text, the document pages were rendered as
images and processed with Tesseract OCR. OCR was required for two
image-based resumes. Extracted text was not manually repaired before
model processing because fragmented words, reading-order errors, and
OCR substitutions represented realistic failures that an ATS-style
pipeline could encounter.

A manual Gold Standard was created for five entity categories:
PERSON, ORG, DATE, GPE, and SKILL. The completed annotation set
contained 1,108 entities across the 20 resumes. PERSON, ORG, DATE, and
GPE represented conventional named-entity categories, while SKILL
captured programming languages, software tools, frameworks, cloud
services, and other technical competencies.

The baseline pipeline applied spaCy's
\texttt{en\_core\_web\_lg} model to identify PERSON, ORG, DATE, and
GPE entities \cite{spacy2020}. A custom keyword matcher was applied
separately to identify SKILL entities, after which the outputs were
combined. Predictions were evaluated against the Gold Standard using
strict exact matching: a prediction was counted as correct only when
its normalized text and entity category matched a Gold Standard entry.
Precision, recall, and F1 were calculated overall, by layout group, and
by entity category.

A second condition evaluated GPT-assisted reconstruction. The same
extracted resume text was supplied to GPT-4o-mini with instructions to
return structured PERSON, ORG, DATE, GPE, and SKILL entities. These
predictions were evaluated independently against the same Gold Standard
and matching rules. This comparison measured whether contextual
reconstruction improved entity recovery over the spaCy and keyword
baseline without changing the underlying extracted text.

Because the layout experiment contained only 20 resumes and was not a
random sample of all resume designs, its results were treated as a
small descriptive comparison rather than evidence of universal
layout effects or proprietary ATS behavior.

\subsection{Large-Scale Resume Corpus}

The second experimental stage expanded the research from the
20-PDF layout study to a larger corpus focused specifically on
technical-skill extraction. The corpus contained 1,020 CS/IT resumes
collected from three public resume datasets. The source datasets were
kept separate so that system performance and prediction distributions
could be reported both by source and in aggregate.

From the Snehaan Bhawal Resume Dataset, 120 resumes labeled for
information-technology-related work were retained
\cite{snehaan_dataset}. These records were assigned identifiers from
\texttt{SNE\_IT\_0001} through \texttt{SNE\_IT\_0120}. From the
Jithin Jagadeesh Resume Dataset, 400 computer-science-oriented records
were retained and assigned identifiers from
\texttt{JIT\_CS\_0001} through \texttt{JIT\_CS\_0400}
\cite{jithin_dataset}. Finally, a strict CS-oriented filter was applied
to the Saugata Roy Arghya Resume Dataset, producing 500 records with
identifiers from \texttt{SAU\_CS\_0001} through
\texttt{SAU\_CS\_0500} \cite{saugata_dataset}.

Each selected record was converted into a consistent JSON-based
representation containing a stable resume identifier, dataset source,
and resume text. A combined metadata file tracked all 1,020 records and
preserved their source labels throughout processing. The resulting
corpus consisted of 500 Saugata records, 400 Jithin records, and 120
Snehaan records.

The source contributions were intentionally unequal because the purpose
of this stage was to evaluate skill extraction across heterogeneous
public datasets rather than construct a perfectly balanced benchmark.
For that reason, subsequent analyses report both source-specific and
combined results. Source-specific reporting helps reveal whether an
overall performance estimate is dominated by one dataset's vocabulary,
writing style, or annotation characteristics.

Unlike the original 20-resume experiment, the large corpus was stored
primarily as extracted resume text rather than consistently preserved
PDF layouts. It was therefore used to evaluate technical-skill
extraction, system scalability, dataset variation, and industry-keyword
coverage. It was not used to draw conclusions about visual resume
layout or reading-order degradation.

\subsection{Technical-Skill Extraction Systems}

Four technical-skill extraction configurations were evaluated on the
large CS/IT corpus:

\begin{enumerate}
    \item \textbf{Generalized V4:} a deterministic, rule-based system
    designed to identify programming languages, software tools,
    frameworks, databases, cloud platforms, and related technical
    competencies. It combined a curated skill vocabulary with
    normalized variants, multi-word phrase matching, and exclusion
    rules intended to reduce common false positives.

    \item \textbf{GPT-only:} a contextual extraction system using
    GPT-4o-mini. Each resume's text was processed independently with
    instructions to return only technical skills in structured JSON
    format. The returned values were normalized and deduplicated within
    each resume before evaluation.

    \item \textbf{Raw hybrid:} the set union of normalized V4 and GPT
    predictions for each resume. This condition retained any skill
    predicted by either system. Every prediction was assigned a
    provenance label indicating whether it appeared in GPT only, V4
    only, or both systems.

    \item \textbf{Filtered hybrid:} a precision-oriented version of the
    raw hybrid. Manually reviewed V4-only predictions were used to learn
    a term-level exclusion filter. The filter was applied only to
    blacklisted V4-only predictions; GPT predictions and terms detected
    by both systems were retained.
\end{enumerate}

A common normalization procedure was applied before comparing or
combining predictions so that differences in capitalization and
surrounding whitespace did not create separate skill entries. The raw
hybrid was intended to increase recall by recovering the complementary
predictions of the rule-based and contextual systems. However, because
a direct union can also introduce rule-based false positives, the
filtered hybrid was designed to recover part of the precision lost by
the raw union.

The exclusion filter was derived from the manually reviewed 50-resume
sample. That review contained 277 V4-only prediction rows representing
66 unique V4-only terms. Seventeen terms that consistently behaved as
false positives were selected for exclusion. Examples of the error
patterns addressed by the filter included generic educational words,
section fragments, and non-technical phrases incorrectly matched by
the rule-based vocabulary.

To avoid evaluating the filter on the same annotations used to learn
it, the 50-resume filtered-hybrid results were generated through
cross-validation. For each held-out fold, exclusion decisions were
derived from the remaining reviewed folds and then applied to the
held-out resumes. After the manual evaluation was completed, a
consolidated 17-term filter learned from the full reviewed sample was
applied to V4-only predictions in the 1,020-resume corpus.

All systems retained the resume identifier and dataset source with each
prediction. This made it possible to compare extraction systems both
across the complete evaluation sample and separately for the Jithin,
Saugata, and Snehaan sources.

\subsection{Manual Evaluation and Metrics}

A stratified sample of 50 resumes was selected from the 1,020-resume
corpus for manual evaluation. The sample contained 15 Jithin resumes,
20 Saugata resumes, and 15 Snehaan resumes. Selection was stratified by
dataset source and initial prediction volume so that resumes with low,
medium, and high numbers of predicted skills were represented.

A review pack was created from the normalized union of V4 and GPT
predictions. It contained 1,338 unique prediction rows, together with
resume-text exports and candidate missed-skill rows. Each prediction
was manually reviewed in the context of its source resume and assigned
one of three primary labels:

\begin{enumerate}
    \item \textbf{Correct:} the predicted term was explicitly present
    and functioned as a technical skill in the resume.

    \item \textbf{Partial:} the prediction captured a valid technical
    concept but did not exactly match the complete skill expression.

    \item \textbf{Incorrect:} the term was absent, non-technical, or
    incorrectly interpreted as a skill.
\end{enumerate}

Skills present in the resume text but not recovered by a particular
system were recorded as false negatives. The completed manual review
therefore provided a shared reference set for evaluating V4, GPT-only,
the raw hybrid, and the cross-validation-filtered hybrid under the same
criteria.

Two scoring schemes were used. Under \emph{strict} scoring, only
predictions labeled correct counted as true positives; partial matches
did not receive full credit. Under \emph{relaxed} scoring, correct and
partial predictions were accepted as successful matches. This
distinction measured whether system rankings changed when near-exact
technical concepts were credited.

For each system, precision, recall, and F1 were calculated as

\begin{equation}
P = \frac{TP}{TP+FP},
\end{equation}

\begin{equation}
R = \frac{TP}{TP+FN},
\end{equation}

and

\begin{equation}
F_1 = \frac{2PR}{P+R},
\end{equation}

where \(TP\), \(FP\), and \(FN\) denote true-positive, false-positive,
and false-negative skill predictions. Results were calculated for the
complete 50-resume sample and separately for the Jithin, Saugata, and
Snehaan sources.

A nonparametric resume-level bootstrap was used to quantify uncertainty
in the strict-F1 difference between the filtered hybrid and GPT-only
systems \cite{efron1993bootstrap}. Resumes were sampled with replacement,
and strict F1 was recomputed for both systems within each bootstrap
sample. The analysis reported the mean F1 difference, a 95\% percentile
confidence interval, and the proportion of bootstrap samples in which
the filtered hybrid achieved a higher F1 score. Resampling was
performed both by dataset source and across the combined evaluation
sample.

Because the manual annotations were completed by one researcher, the
evaluation did not provide an inter-annotator agreement estimate. This
constraint is addressed as a threat to validity later in the report.

\subsection{External Comparison Methods}

Two exploratory external analyses were conducted to connect the
evaluated extraction results with current job-market vocabulary and a
commercial resume-scanning service.

\textit{Industry-keyword comparison:}
Fifty Computer Science and Software Engineering job descriptions were
collected to provide a sample of technical terminology appearing in
current job advertisements. Technical keywords from these descriptions
were compiled and their frequencies were summarized. The resulting
industry vocabulary was compared with the manually annotated SKILL
entities and system predictions from the original 20-resume experiment.

The comparison evaluated whether industry-relevant skills present in the
Gold Standard were recovered by the original baseline or the
GPT-assisted system. The baseline contributed 207 deduplicated SKILL
predictions. Exact and normalized matching rules were applied
consistently when comparing the industry vocabulary, Gold Standard
skills, and system outputs. Results were summarized overall and for the
standard and creative or complex resume groups.

This analysis was intended to provide descriptive job-market context
rather than a representative estimate of all CS/IT hiring requirements.
The 50 job descriptions were not treated as a new manually annotated
resume benchmark, and keyword frequency was not interpreted as evidence
of proprietary ATS ranking behavior.

\textit{Jobscan pilot:}
A small external pilot was conducted using Jobscan, a proprietary resume
scanner and job-match reporting service \cite{jobscan2026}. Four valid
resume--job-description scans were retained: one Jithin resume
(\texttt{JIT\_CS\_0016}) and three Snehaan resumes
(\texttt{SNE\_IT\_0053}, \texttt{SNE\_IT\_0056}, and
\texttt{SNE\_IT\_0110}). Each resume was paired with a different
role-relevant job description.

For each valid scan, the recorded outputs included the overall Match
Score, visible matched and missing skills, hard- and soft-skill
observations, searchability feedback, recruiter recommendations, and
word-count information. An initial scan of
\texttt{SNE\_IT\_0053} returned an anomalous score of 100 after the wrong
text was likely submitted. That result was excluded, the scan was
repeated with the intended content, and only the corrected result was
included in the pilot summary.

Jobscan results were analyzed descriptively and were not converted into
NER precision, recall, or F1. Its Match Score incorporates factors beyond
technical-skill extraction, including headings, contact information,
education, job titles, measurable results, word count, and general role
alignment. In addition, the reports did not expose every possible missing
skill and occasionally produced questionable role or education
interpretations. Jobscan was therefore treated as an exploratory
industry-facing comparison rather than an external Gold Standard or a
verified representation of all commercial ATS platforms.

\section{Results}

\subsection{Original Layout Experiment}

The original baseline produced 1,277 predictions for the 1,108
Gold Standard entities. Of these predictions, 322 were correct and 955
were incorrect, while 786 Gold Standard entities were missed. The
resulting combined strict precision, recall, and F1 were 0.252, 0.291,
and 0.270, respectively.

Table~\ref{tab:layout_system_results} compares the baseline and
GPT-assisted reconstruction results by layout group. The baseline
achieved a strict F1 of 0.223 on the 10 standard resumes and 0.329 on
the 10 creative or complex resumes. The creative group therefore
outperformed the standard group by 0.106 F1, contrary to the original
directional hypothesis.

\begin{table}[t]
\caption{Strict Performance in the Original 20-Resume Experiment}
\label{tab:layout_system_results}
\centering
\footnotesize
\begin{tabular}{llccc}
\toprule
Group & System & Precision & Recall & F1 \\
\midrule
Standard & Baseline & 0.211 & 0.236 & 0.223 \\
Standard & GPT & 0.599 & 0.273 & 0.375 \\
Creative & Baseline & 0.302 & 0.360 & 0.329 \\
Creative & GPT & 0.708 & 0.424 & 0.530 \\
Combined & Baseline & 0.252 & 0.291 & 0.270 \\
Combined & GPT & 0.654 & 0.339 & 0.447 \\
\bottomrule
\end{tabular}
\end{table}

Baseline performance varied substantially across entity categories, as
shown in Table~\ref{tab:baseline_entity_results}. ORG produced the
lowest F1, 0.125, because the general-purpose NER model generated 788
organization predictions but only 56 were correct. PERSON achieved high
recall of 0.778 but low precision of 0.206. The SKILL matcher showed the
opposite pattern: it achieved the highest precision, 0.894, but recovered
only 0.235 of the 788 annotated skills.

\begin{table}[t]
\caption{Baseline Strict Performance by Entity Category}
\label{tab:baseline_entity_results}
\centering
\scriptsize
\setlength{\tabcolsep}{3.5pt}
\begin{tabular}{lrrrrrr}
\toprule
Entity & Gold & Pred. & Correct & P & R & F1 \\
\midrule
PERSON & 18 & 68 & 14 & 0.206 & 0.778 & 0.326 \\
ORG & 111 & 788 & 56 & 0.071 & 0.505 & 0.125 \\
DATE & 146 & 142 & 47 & 0.331 & 0.322 & 0.326 \\
GPE & 45 & 72 & 20 & 0.278 & 0.444 & 0.342 \\
SKILL & 788 & 207 & 185 & 0.894 & 0.235 & 0.372 \\
\bottomrule
\end{tabular}
\end{table}

GPT-assisted reconstruction generated 575 predictions, of which 376
were correct and 199 were incorrect; 732 Gold Standard entities remained
unrecovered. Its combined strict precision, recall, and F1 were 0.654,
0.339, and 0.447. Relative to the baseline, GPT increased strict F1 by
0.177 overall. Improvements occurred in both layout groups: standard
resume F1 increased from 0.223 to 0.375, while creative or complex resume
F1 increased from 0.329 to 0.530.

These findings indicate that contextual reconstruction improved entity
recovery from the same extracted text. However, the higher performance
of the creative group does not demonstrate that creative layouts are
generally easier to parse. The result may reflect differences in
profession, entity density, skill vocabulary, anonymization, OCR use,
and document composition within the small 20-resume sample.

\subsection{Fifty-Resume System Comparison}

The completed 50-resume manual review contained 1,592 Gold Standard
technical skills. Generalized V4 generated 488 predictions: 287 were
correct, 189 were incorrect, and 12 were partial matches. GPT-only
generated 1,061 predictions, including 1,025 correct, 24 incorrect, and
12 partial predictions. The raw hybrid union generated 1,338 predictions,
of which 1,104 were correct, 211 were incorrect, and 23 were partial.

The cross-validation-filtered hybrid produced 1,174 predictions,
including 1,081 correct, 79 incorrect, and 14 partial predictions.
Its reported performance was generated by applying exclusion decisions
learned from training folds to held-out resumes rather than by evaluating
the full-sample filter on the annotations used to construct it.

Table~\ref{tab:system_comparison_50} presents the combined strict and
relaxed evaluation results. Generalized V4 achieved strict precision of
0.588 but recovered only 0.180 of the Gold Standard skills, producing a
strict F1 of 0.276. GPT-only substantially improved both precision and
recall, reaching a strict F1 of 0.773.

\begin{table}[t]
\caption{Combined 50-Resume Technical-Skill Evaluation}
\label{tab:system_comparison_50}
\centering
\scriptsize
\setlength{\tabcolsep}{2.0pt}
\begin{tabular}{lrrrrrrr}
\toprule
& & \multicolumn{3}{c}{Strict} &
\multicolumn{3}{c}{Relaxed} \\
\cmidrule(lr){3-5}\cmidrule(lr){6-8}
System & Pred. & P & R & F1 & P & R & F1 \\
\midrule
V4 & 488 & 0.588 & 0.180 & 0.276 &
0.613 & 0.188 & 0.288 \\
GPT-only & 1,061 & 0.966 & 0.644 & 0.773 &
0.977 & 0.651 & 0.782 \\
Raw hybrid & 1,338 & 0.825 & 0.693 & 0.754 &
0.842 & 0.708 & 0.769 \\
Filtered hybrid (CV) & 1,174 & 0.921 & 0.679 & 0.782 &
0.933 & 0.688 & 0.792 \\
\bottomrule
\end{tabular}
\end{table}

The raw hybrid achieved the highest strict recall, 0.693, because it
retained every normalized skill predicted by either V4 or GPT. However,
its strict precision declined to 0.825 because V4-only false positives
were also added to the union. Consequently, its strict F1 of 0.754 was
lower than the GPT-only F1 of 0.773.

Cross-validation filtering removed 164 V4-only predictions from the raw
union. These removals included 132 incorrect predictions and 23 correct
predictions. The resulting filtered hybrid recovered most of the raw
hybrid's precision loss while retaining higher recall than GPT-only. It
achieved strict precision of 0.921, recall of 0.679, and the highest
combined strict F1, 0.782.

Under relaxed scoring, the system ordering remained similar. The filtered
hybrid achieved the highest relaxed F1, 0.792, followed by GPT-only at
0.782, the raw hybrid at 0.769, and V4 at 0.288. The filtered hybrid's
combined strict-F1 advantage over GPT-only was 0.009, indicating a modest
overall improvement despite the larger reduction in rule-based false
positives.

\subsection{Dataset-Specific Performance}

Performance differed across the three dataset sources. The manually
reviewed sample contained 274 Gold Standard skills in the 15 Jithin
resumes, 577 in the 20 Saugata resumes, and 741 in the 15 Snehaan
resumes. Thus, the equally sized Jithin and Snehaan groups differed
substantially in technical-skill density.

Table~\ref{tab:dataset_system_results} reports strict performance by
source. GPT-only achieved the highest strict F1 for Jithin, 0.827, and
Saugata, 0.841. The cross-validation-filtered hybrid was close behind,
with strict F1 scores of 0.822 and 0.833, respectively.

\begin{table}[t]
\caption{Strict Performance by Dataset Source}
\label{tab:dataset_system_results}
\centering
\scriptsize
\setlength{\tabcolsep}{3.0pt}
\begin{tabular}{llccc}
\toprule
Dataset & System & Precision & Recall & F1 \\
\midrule
Jithin & V4 & 0.678 & 0.223 & 0.335 \\
Jithin & GPT-only & 0.966 & 0.723 & \textbf{0.827} \\
Jithin & Raw hybrid & 0.854 & 0.745 & 0.795 \\
Jithin & Filtered hybrid & 0.923 & 0.741 & 0.822 \\
\midrule
Saugata & V4 & 0.593 & 0.194 & 0.292 \\
Saugata & GPT-only & 0.973 & 0.740 & \textbf{0.841} \\
Saugata & Raw hybrid & 0.835 & 0.769 & 0.801 \\
Saugata & Filtered hybrid & 0.922 & 0.759 & 0.833 \\
\midrule
Snehaan & V4 & 0.545 & 0.154 & 0.240 \\
Snehaan & GPT-only & 0.959 & 0.540 & 0.691 \\
Snehaan & Raw hybrid & 0.804 & 0.615 & 0.697 \\
Snehaan & Filtered hybrid & 0.919 & 0.594 & \textbf{0.721} \\
\bottomrule
\end{tabular}
\end{table}

Snehaan showed a different pattern. GPT-only achieved strict precision
of 0.959 but recall of only 0.540, producing an F1 of 0.691. The raw
hybrid increased recall to 0.615 but reduced precision to 0.804. After
cross-validation filtering, precision recovered to 0.919 while recall
remained above GPT-only at 0.594. The filtered hybrid consequently
achieved the highest Snehaan strict F1, 0.721, an improvement of 0.030
over GPT-only.

The raw hybrid achieved the highest strict recall within every source,
but its added V4-only predictions consistently reduced precision.
Filtering produced a more balanced tradeoff. On Jithin and Saugata, this
tradeoff was not sufficient to exceed GPT-only F1; on Snehaan, the
additional recovered skills outweighed the remaining false positives.

Relaxed scoring did not change the winning GPT-based system within each
source. GPT-only achieved relaxed F1 scores of 0.852 on Jithin and 0.841
on Saugata, while the filtered hybrid achieved the highest Snehaan
relaxed F1, 0.733.

Generalized V4 performed substantially below the GPT-based systems for
all three sources, with strict F1 ranging from 0.240 to 0.335. Its low
recall indicates that a fixed rule-based vocabulary did not cover much
of the varied technical terminology present in the manually reviewed
resumes.

These differences demonstrate why combined results alone are
insufficient. The filtered hybrid's overall advantage was driven
primarily by its improvement on Snehaan, while GPT-only remained
slightly stronger on Jithin and Saugata. Because these sources contain
different text structures, vocabularies, and skill densities, the
dataset-specific differences should not be interpreted as visual-layout
effects.

\subsection{Bootstrap Uncertainty}

Resume-level bootstrap resampling was used to evaluate whether the
strict-F1 differences between the cross-validation-filtered hybrid and
GPT-only systems were stable across alternative samples of resumes.
Positive differences favor the filtered hybrid, while negative
differences favor GPT-only.

Table~\ref{tab:bootstrap_results} reports the mean F1 differences,
95\% percentile confidence intervals, and the proportion of bootstrap
samples in which the filtered hybrid achieved a higher strict F1.

\begin{table}[t]
\caption{Bootstrap Strict-F1 Difference: Filtered Hybrid Minus GPT-Only}
\label{tab:bootstrap_results}
\centering
\scriptsize
\setlength{\tabcolsep}{3.0pt}
\begin{tabular}{lccc}
\toprule
Group & $\Delta$F1 & 95\% CI & $P(\Delta > 0)$ \\
\midrule
Combined & $+0.0089$ & $[-0.0002,\ 0.0179]$ & 0.9725 \\
Jithin & $-0.0049$ & $[-0.0186,\ 0.0075]$ & 0.2362 \\
Saugata & $-0.0079$ & $[-0.0249,\ 0.0078]$ & 0.1682 \\
Snehaan & $+0.0305$ & $[0.0158,\ 0.0457]$ & 1.0000 \\
\bottomrule
\end{tabular}
\end{table}

For the combined 50-resume sample, the filtered hybrid improved strict
F1 by 0.0089 on average. It outperformed GPT-only in 97.25\% of the
bootstrap samples. However, the 95\% confidence interval ranged from
$-0.0002$ to $0.0179$ and therefore slightly crossed zero. The combined
result provides strong directional evidence of improvement, but it does
not establish an unambiguous positive difference under the 95\%
interval criterion.

The clearest improvement occurred for Snehaan. The filtered hybrid's
mean strict-F1 advantage was 0.0305, with a 95\% confidence interval
from 0.0158 to 0.0457. The filtered hybrid outperformed GPT-only in
every Snehaan bootstrap sample. This interval remained entirely above
zero, indicating that the source-specific improvement was stable under
resume-level resampling.

Jithin and Saugata showed small negative point estimates of
$-0.0049$ and $-0.0079$, respectively. Both confidence intervals
included zero, so the analysis did not establish a stable difference
between GPT-only and the filtered hybrid for either source.

These results confirm that the filtered hybrid's effectiveness was
dataset-dependent. It produced a clear gain for Snehaan, while GPT-only
remained slightly stronger on Jithin and Saugata. The quantity
$P(\Delta > 0)$ is a descriptive bootstrap proportion and should not
be interpreted as a conventional hypothesis-test p-value.

\subsection{Full-Corpus Skill Recovery}

After completing the manually reviewed evaluation, the extraction systems
were applied to all 1,020 resumes. This stage measured prediction volume,
source distribution, and scalability rather than accuracy, because the
complete corpus did not have manually annotated Gold Standard skills.

The Generalized V4 output initially contained 14,665 entity predictions.
After removing 5,242 non-SKILL entities, 9,423 technical-skill
predictions remained. V4 produced at least one retained skill for 1,019
of the 1,020 resumes.

GPT-only processing completed successfully for all 1,020 resumes and
generated 20,756 normalized skill predictions. Of the complete corpus,
1,016 resumes contained at least one GPT-predicted skill.

The normalized raw union contained 25,238 unique resume-skill
predictions. Of these, 15,815 were predicted only by GPT, 4,482 were
predicted only by V4, and 4,941 were detected by both systems. The raw
hybrid total was therefore smaller than the direct sum of the two system
outputs because shared predictions were deduplicated.

The consolidated 17-term exclusion filter was applied only to the 4,482
V4-only predictions. It removed 2,502 predictions and retained 1,980
V4-only predictions. All 15,815 GPT-only predictions and all 4,941
predictions detected by both systems were retained. The resulting
filtered-hybrid corpus contained 22,736 predictions, a 9.9\% reduction
from the raw union.

Table~\ref{tab:full_corpus_counts} reports the effect of filtering by
dataset source.

\begin{table}[t]
\caption{Full-Corpus Raw and Filtered Hybrid Prediction Counts}
\label{tab:full_corpus_counts}
\centering
\footnotesize
\setlength{\tabcolsep}{3.0pt}
\begin{tabular}{lrrrr}
\toprule
Dataset & Resumes & Raw & Filtered & Removed \\
\midrule
Jithin & 400 & 6,458 & 5,994 & 464 \\
Saugata & 500 & 13,546 & 12,157 & 1,389 \\
Snehaan & 120 & 5,234 & 4,585 & 649 \\
\midrule
Combined & 1,020 & 25,238 & 22,736 & 2,502 \\
\bottomrule
\end{tabular}
\end{table}

The proportion removed differed by source. Filtering removed
approximately 7.2\% of Jithin predictions, 10.3\% of Saugata
predictions, and 12.4\% of Snehaan predictions. Snehaan therefore showed
the largest proportional change, consistent with the manually reviewed
sample in which filtering produced the clearest source-specific benefit.

These full-corpus totals should not be interpreted as precision, recall,
or F1 results. Prediction volume can reflect resume length, technical
skill density, vocabulary, source formatting, and false-positive
behavior. Accuracy conclusions remain limited to the manually reviewed
50-resume evaluation, while the 1,020-resume analysis demonstrates the
scale and distributional effect of applying the extraction and filtering
pipeline.

\subsection{External and Industry Comparisons}

\textit{Industry-keyword recovery:}
The 50-job-description analysis produced a vocabulary of 496 technical
keywords. Of the 788 Gold Standard SKILL occurrences in the original
20 resumes, 190 matched this broader industry vocabulary. These included
96 occurrences in the standard group and 94 in the creative or complex
group.

Under normalized matching, the original baseline recovered 105 of the
190 industry-relevant skill occurrences, or 55.3\%. GPT recovered 99,
or 52.1\%. The baseline therefore achieved a 3.2-percentage-point
advantage across the complete 496-keyword vocabulary.

The group-level results were not uniform. On standard resumes, GPT
recovered 44 of 96 occurrences, or 45.8\%, compared with 43, or 44.8\%,
for the baseline. On creative or complex resumes, the baseline recovered
62 of 94 occurrences, or 66.0\%, while GPT recovered 55, or 58.5\%.

A second analysis focused on the 18 leading keywords that appeared in
at least seven of the collected job descriptions. Forty-four occurrences
of these leading keywords appeared in the Gold Standard resumes. The
baseline recovered 32 occurrences, or 72.7\%, while GPT recovered 33,
or 75.0\%.

Table~\ref{tab:industry_recovery} summarizes the normalized recovery
results. GPT performed slightly better for leading keywords overall and
for the standard group, while the baseline remained stronger across the
complete vocabulary and the creative group.

\begin{table}[t]
\caption{Normalized Industry-Keyword Recovery}
\label{tab:industry_recovery}
\centering
\scriptsize
\setlength{\tabcolsep}{2.5pt}
\begin{tabular}{llrrr}
\toprule
Scope & Group & Occurrences & Baseline & GPT \\
\midrule
All 496 & Standard & 96 & 44.8\% & 45.8\% \\
All 496 & Creative & 94 & 66.0\% & 58.5\% \\
All 496 & Combined & 190 & 55.3\% & 52.1\% \\
\midrule
Leading 18 & Standard & 25 & 60.0\% & 64.0\% \\
Leading 18 & Creative & 19 & 89.5\% & 89.5\% \\
Leading 18 & Combined & 44 & 72.7\% & 75.0\% \\
\bottomrule
\end{tabular}
\end{table}

Under exact matching, the baseline recovered 102 of 190 occurrences
across the full vocabulary, compared with 83 for GPT. For the leading
keywords, the baseline recovered 30 of 44 occurrences and GPT recovered
27. The difference between exact and normalized recovery demonstrates
the effect of capitalization, punctuation, abbreviations, and related
surface-form variation.

\textit{Jobscan pilot:}
Four valid resume--job-description scans were retained for the
exploratory Jobscan comparison. Their Match Scores are shown in
Table~\ref{tab:jobscan_results}.

\begin{table}[t]
\caption{Valid Jobscan Pilot Scores}
\label{tab:jobscan_results}
\centering
\footnotesize
\begin{tabular}{llr}
\toprule
Resume & Source & Match Score \\
\midrule
\texttt{JIT\_CS\_0016} & Jithin & 30 \\
\texttt{SNE\_IT\_0053} & Snehaan & 27 \\
\texttt{SNE\_IT\_0056} & Snehaan & 50 \\
\texttt{SNE\_IT\_0110} & Snehaan & 21 \\
\bottomrule
\end{tabular}
\end{table}

The four valid scores had a mean of 32.0, a median of 28.5, and a range
from 21 to 50. Across the visible report fields, Jobscan identified 11
matched skills and 33 missing skills. The initial
\texttt{SNE\_IT\_0053} scan that returned a score of 100 was excluded
because the wrong text was likely submitted; the corrected scan produced
the retained score of 27.

These Jobscan outcomes were not directly comparable with the study's
precision, recall, or F1 results. Each resume was evaluated against a
different role-specific job description, and the Match Score incorporated
formatting, searchability, education, job titles, measurable results,
word count, and general role alignment in addition to visible skills.
The four-scan pilot therefore supplied limited external context rather
than independent validation of the extraction systems.

Together, the external comparisons did not identify one system as
universally superior. GPT showed a small advantage for leading
industry keywords, the baseline recovered more of the complete industry
vocabulary, and Jobscan produced substantial variation across only four
resume--job pairs.

\section{Discussion}

\subsection{Layout Hypothesis}

The original directional hypothesis predicted that creative or complex
resumes would achieve lower strict entity-level F1 than standard
resumes. The results did not support this hypothesis. Under the original
baseline, the creative group achieved an F1 of 0.329, compared with
0.223 for the standard group. The same ordering remained after
GPT-assisted reconstruction: creative or complex resumes achieved 0.530,
while standard resumes achieved 0.375.

The entity-specific results show that the difference was not consistent
across every category. Standard resumes achieved higher DATE F1
(0.373 versus 0.294) and GPE F1 (0.407 versus 0.286). Creative resumes
achieved higher ORG F1 (0.188 versus 0.082), PERSON F1
(0.444 versus 0.000), and SKILL F1 (0.456 versus 0.313). The aggregate
layout-group difference therefore reflected the mixture and distribution
of entity categories rather than a uniform advantage for creative
formatting.

This result should not be interpreted as evidence that creative layouts
are inherently easier for automated systems. Layout type was not the
only characteristic that differed between the two groups. The resumes
also varied in profession, technical-skill density, organization usage,
anonymization, wording, document length, and extraction method. Some
creative resumes contained explicit lists of technical skills that were
well suited to keyword matching, while some standard resumes contained
organization-heavy text that generated many spaCy false positives.

The sample also contained only 10 resumes in each layout group. The
standard and creative resumes were not paired versions of identical
content, so the experiment could not isolate layout while holding
vocabulary and entity density constant. OCR was required for two
image-based documents, and strict exact matching penalized boundary,
capitalization, abbreviation, and normalization differences. These
factors may have affected group-level results independently of visual
design.

Resume parsing involves at least two related stages: converting the PDF
into correctly ordered text and identifying entities within that text.
A visually creative document can perform adequately when its embedded
text and reading order remain accessible, while a visually conventional
document can still produce poor predictions if its content is
anonymized, organization-heavy, or outside the system vocabulary.
Layout-aware research similarly emphasizes that document heterogeneity
and reading-order reconstruction must be evaluated separately from
semantic extraction \cite{zhu2025layout}.

Accordingly, the most defensible conclusion is not that creative layouts
are harmless, but that this small experiment did not establish a
universal creative-layout penalty. A stronger causal test would place
the same resume content into multiple controlled templates, preserve
identical Gold Standard entities, and evaluate PDF extraction and NER as
separate stages. Such a design would allow differences to be attributed
more directly to layout rather than to resume content or dataset
composition.

For practical guidance, the findings support caution rather than a
complete prohibition on visual design. Important information should
remain selectable, headings should be recognizable, reading order should
be testable after text extraction, and essential qualifications should
not appear only inside graphics or images.

\subsection{Hybrid-System Tradeoffs}

The four technical-skill extraction configurations represented different
precision, recall, interpretability, and implementation tradeoffs.
Generalized V4 was deterministic and auditable, but its strict recall of
0.180 showed that a fixed vocabulary and pattern rules could not cover
the terminology present across the three datasets. Its strict precision
of 0.588 also demonstrated that rule-based matching did not automatically
guarantee high precision when generic or context-dependent terms appeared
in resume text.

GPT-only provided the strongest standalone performance. Its strict
precision of 0.966 indicates that contextual extraction rejected most
non-technical phrases, while its recall of 0.644 recovered substantially
more skills than V4. This result suggests that contextual interpretation
was particularly valuable for abbreviations, multi-word tools, uncommon
technologies, and terms whose meaning depended on surrounding resume
content.

The raw hybrid confirmed that the systems contained complementary
predictions, but it also demonstrated that set union is not automatically
an improvement. Adding every V4 prediction to GPT increased strict
recall from 0.644 to 0.693. The raw union recovered 79 more correct
predictions than GPT-only, but it also introduced 187 additional
incorrect predictions. Precision therefore declined from 0.966 to 0.825,
and strict F1 fell from 0.773 to 0.754.

The cross-validation-filtered hybrid produced a better balance. Filtering
removed 164 V4-only predictions from the evaluated union, including 132
incorrect, 23 correct, and nine partial predictions. Strict precision
recovered to 0.921 while recall remained above GPT-only at 0.679. The
resulting strict F1 of 0.782 was the highest combined score, and the
relaxed F1 of 0.792 showed the same overall pattern.

However, filtering did not improve every source equally. GPT-only
retained the highest strict F1 for Jithin and Saugata, while the filtered
hybrid performed best for Snehaan. The combined bootstrap difference
favored the filtered hybrid by 0.0089 F1, but its 95\% confidence interval
slightly crossed zero. In contrast, the Snehaan improvement of 0.0305
had an interval entirely above zero. These findings indicate that the
value of rule-based additions depended on dataset vocabulary and skill
density rather than representing a universal hybrid advantage.

The full-corpus application introduces an additional transfer concern.
The consolidated filter removed 2,502 of the 4,482 V4-only predictions
from the raw 1,020-resume union. Although the reviewed sample showed
that many excluded terms behaved as false positives, the full corpus
did not contain complete Gold Standard annotations. It is therefore not
possible to assume that every removed full-corpus prediction was
incorrect. A term-level blacklist learned from one sample may remove a
valid skill when the same expression appears in a different technical
context.

The systems also differ operationally. V4 can run locally at low cost,
produces repeatable outputs, and permits direct inspection of every rule.
GPT provides stronger contextual coverage but depends on model access,
prompt stability, API cost, and privacy-aware handling of resume text.
The filtered hybrid adds further maintenance because it requires
prediction provenance, reviewed examples, exclusion rules, and held-out
validation.

Accordingly, GPT-only represents a strong and comparatively simple
default for this corpus. A filtered hybrid can provide additional recall
and modestly higher combined F1 when its exclusion rules are validated
for the target source. The raw union should not be preferred solely
because it returns more skills. Practical hybrid design requires
provenance tracking, source-specific evaluation, and explicit control of
the false positives introduced by each component.

\subsection{Practical Implications}

For students and resume writers, the findings do not support a simple
rule that every visually creative resume will be parsed less accurately
than a standard resume. Instead, they emphasize the importance of the
machine-readable structure beneath the visual design. Important resume
content should remain selectable, copyable, and understandable after
plain-text extraction. Applicants can perform a basic self-audit by
copying the PDF content into a text editor and checking whether names,
contact details, education, employers, dates, locations, skills, and
section headings appear in a logical reading order. Essential
qualifications should not be presented only through icons, images,
decorative text boxes, or graphics.

Technical skills should be written explicitly and truthfully. Multi-word
technology names should remain intact, and uncommon abbreviations may be
paired with their expanded forms when doing so improves clarity. Relevant
skills can be placed in a dedicated skills section and reinforced through
accurate project or experience descriptions. However, applicants should
not add technologies they have not used or repeat keywords merely to
increase their frequency. The objective is to make genuine qualifications
clear to both human readers and automated systems.

The job-description experiments also indicate that keyword comparison
depends on the extraction method. Across all 496 normalized job-description
terms, the deterministic baseline recovered 55.3\%, compared with 52.1\%
for GPT. Among the 18 leading industry keywords, GPT recovered 75.0\%,
compared with 72.7\% for the baseline. Neither system therefore dominated
every comparison. A useful resume review should combine multiple checks:
confirm that relevant terminology is present when accurate, inspect the
surrounding context, and manually review any terms reported as missing.

The four valid Jobscan pilot scores ranged from 21 to 50, with a mean of
32. These scores should be interpreted as exploratory system outputs
rather than direct measurements of resume quality or hiring probability.
Commercial matching tools may incorporate formatting, section structure,
experience wording, job-title similarity, and other proprietary factors
in addition to technical skills. Applicants should therefore avoid
optimizing a resume around one score. Automated feedback is more useful
as a prompt for human review than as a definitive pass-or-fail judgment.

For developers and researchers, the results demonstrate the importance
of reporting precision and recall together. The raw hybrid increased
recall but also introduced enough false positives to reduce F1 below the
GPT-only result. The filtered hybrid achieved the highest combined F1,
but its improvement varied by dataset source and depended on
cross-validation-derived exclusion rules. Hybrid systems should therefore
retain prediction provenance, validate filtering rules on held-out data,
and report source-specific results rather than relying only on an
aggregate score. Terms removed from the unannotated full corpus also
cannot automatically be treated as errors.

Any resume-processing workflow should additionally consider privacy.
Resume text can contain names, contact information, employment history,
education, and other sensitive personal data. When external models or
commercial services are used, researchers and applicants should consider
what information is transmitted, how it is stored, and whether identifying
details can be removed before processing.

Consequently, the student-facing self-audit rubric and optimization guide
developed from this study should present conservative, evidence-based
checks rather than promises of ATS success. The guidance can help users
identify machine-readability problems, missing relevant terminology, and
potential extraction failures, but it cannot reproduce every commercial
ATS or predict a hiring outcome. Automated analysis should complement
careful human editing and review.

\section{Limitations and Threats to Validity}

Several limitations constrain the interpretation and generalizability of
this study. First, the layout experiment used only 20 resumes, divided
equally between standard and creative or complex designs. The documents
were not paired versions of identical resume content. They differed in
profession, vocabulary, entity density, document length, anonymization,
and use of OCR. Consequently, the observed performance differences cannot
be attributed solely to layout. A stronger design would place identical
content into multiple controlled templates and evaluate text extraction
and entity recognition separately.

The standard and creative categories also required human judgment.
Although the categories captured visible structural differences, resume
design exists on a continuum and cannot always be represented by a
binary label. Features such as columns, icons, tables, graphics, embedded
fonts, and reading order may affect extraction differently. The current
sample was too small to estimate the independent effect of each feature.

The Gold Standard annotations introduce another source of uncertainty.
Annotations were manually prepared using a predefined entity schema, but
the study did not include multiple independent annotators or calculate
inter-annotator agreement. Some organization names, locations, dates, and
technical skills can reasonably be represented using different boundaries
or levels of specificity. Strict exact matching counted capitalization,
abbreviation, normalization, and boundary differences as errors even when
a prediction retained useful information. Relaxed evaluation reduced this
problem for the technical-skill experiment, but it still depended on
manual review decisions.

The large-scale corpus was assembled from three public resume datasets
and restricted to computer-science-related records. These sources differed
in format, vocabulary, category definitions, anonymization, and technical
skill density. The 1,020 resumes should therefore not be interpreted as a
representative sample of all applicants, professions, geographic regions,
or commercial ATS inputs. Public datasets may also contain outdated,
duplicated, synthetic, or heavily processed records that differ from
resumes submitted in current hiring environments.

The manually reviewed 50-resume sample was stratified by dataset source
and prediction-count band. This design ensured that low-, medium-, and
high-output cases were inspected, but it was not a simple random sample
of the full corpus. The resulting precision, recall, and F1 estimates are
appropriate for comparing the evaluated systems within the sample, but
they may not directly estimate performance across all 1,020 resumes.

Although cross-validation was used when evaluating the filtered hybrid,
the consolidated exclusion terms were derived from reviewed V4-only
predictions within the available datasets. Their effectiveness may depend
on the vocabulary and skill distribution of those sources. The filtered
hybrid improved combined strict F1 by approximately 0.009 over GPT-only,
but the bootstrap confidence interval slightly crossed zero. The
improvement was clearly positive for Snehaan but was not reproduced for
Jithin or Saugata. Independent evaluation on a new annotated dataset is
therefore necessary before treating the filter as generally transferable.

The full 1,020-resume corpus did not receive complete Gold Standard
annotation. As a result, the 22,736 filtered-hybrid predictions describe
system output volume rather than verified accuracy. Similarly, the 2,502
V4-only predictions removed by the consolidated filter cannot all be
assumed to be incorrect. Some excluded expressions may represent valid
skills in contexts not observed during manual review.

GPT-based results are also dependent on the selected model, prompt,
normalization procedure, API behavior, and time of execution. Model
updates or alternative prompts may produce different predictions.
Although structured outputs and a consistent processing pipeline improved
repeatability, the system is not deterministic in the same manner as the
rule-based baseline.

Finally, this research evaluated an experimental pipeline rather than a
complete commercial ATS. The Jobscan comparison contained only four valid
resume--job-description pairs, with different job descriptions used for
each resume. Its scores were exploratory and were not directly comparable
with entity-level precision, recall, or F1. The study therefore supports
conclusions about the evaluated extraction methods, not claims about every
commercial screening system or a candidate's probability of receiving an
interview.

\section{Privacy, Ethics, and Reproducibility}

Resumes contain potentially sensitive personal information, including
names, contact details, education, employment history, locations, and
professional affiliations. Although this research used resumes obtained
from publicly accessible datasets rather than directly recruiting
participants, public availability does not eliminate the need for careful
data handling. Raw resume files, extracted text, detailed annotations,
model outputs, and manual-review materials were therefore treated as
research data rather than as information intended for redistribution.

Project files containing resume-level text or predictions were kept
separate from aggregate research results. Private output directories were
excluded from version-controlled research materials, while shareable
outputs were limited to scripts, summary statistics, aggregate tables,
and non-identifying methodological documentation. API credentials were
supplied through environment variables and were not written into scripts
or committed to the repository. Any future public release should continue
to exclude raw personal information unless the relevant dataset terms and
ethical review explicitly permit redistribution.

The use of an external language-model API introduces an additional
privacy consideration because resume text may be transmitted outside the
local research environment. Researchers should minimize the amount of
personal information submitted, remove unnecessary identifying details
when possible, and review the provider's data-use and retention terms.
Similar caution applies to third-party resume-scoring services such as
Jobscan. The limited Jobscan pilot was used only as an exploratory
comparison and should not be interpreted as authorization to upload
another person's resume without their knowledge.

This study was designed to evaluate information-extraction methods rather
than rank candidates or make employment decisions. The systems were not
used to infer protected characteristics, recommend applicants, or predict
job performance. Nevertheless, extraction errors can have ethical
consequences if such systems are incorporated into hiring workflows.
False negatives may hide legitimate qualifications, while false positives
may misrepresent an applicant's experience. Differences across dataset
sources further demonstrate that a system that performs well for one
group of documents may not perform equally well for another.

For these reasons, automated resume analysis should include human review,
especially when its output could influence access to employment.
Applicants should also not be encouraged to believe that a particular
format, keyword count, or commercial score guarantees favorable
treatment. The student-facing guidance produced by this project is
intended to improve clarity and machine readability while preserving
truthful representation of qualifications.

Reproducibility was supported through a structured project directory,
separate scripts for dataset preparation and each extraction system,
normalized CSV schemas, documented prompts, source-specific results, and
saved aggregate summaries. The deterministic V4 pipeline can be rerun
using the same vocabulary, normalization rules, and exclusion logic.
The filtered-hybrid evaluation also recorded its cross-validation
procedure, while the statistical analysis documented resume-level
bootstrap resampling and its reported confidence intervals.

GPT-based processing is less completely reproducible because outputs may
change with model versions, service updates, prompts, and nondeterministic
generation. Reproduction therefore requires recording the exact model
identifier, prompt structure, processing date, normalization procedure,
and output schema. Saved status files and combined prediction outputs
provide an audit trail for the completed runs, but they cannot guarantee
that a future API execution will return identical predictions.

Finally, reproducibility should not require releasing sensitive
resume-level data. Researchers can provide code, configuration details,
entity schemas, aggregate evaluation results, and instructions for
rebuilding the dataset from its original sources while keeping private
annotations and identifying content restricted. This approach balances
scientific transparency with the privacy risks associated with resume
data.

\section{Student-Facing Deliverables}

The experimental findings were translated into two practical
student-facing deliverables: a Resume Self-Audit Rubric and a Computer
Science Resume Optimization Guide. These resources emphasize
machine-readable structure, truthful skill representation, and manual
verification. They are not intended to reproduce a commercial ATS,
predict hiring outcomes, or guarantee that a resume will pass an
automated screening system.

\subsection{Resume Self-Audit Rubric}

The Resume Self-Audit Rubric provides a structured way for students to
inspect a resume before submission. Each criterion may be rated as zero
when it is not satisfied, one when it is partially satisfied, or two when
it is clearly satisfied. The total is intended only to identify areas
requiring revision; it is not a validated ATS compatibility score.

\begin{enumerate}
    \item \textbf{Text accessibility:} The PDF contains selectable and
    copyable text. Important content is not stored only as an image.

    \item \textbf{Reading order:} When the resume is copied into a plain
    text editor, sections, bullet points, dates, and descriptions appear
    in a logical sequence.

    \item \textbf{Section recognition:} Conventional and descriptive
    headings such as Education, Experience, Projects, and Technical
    Skills are used where appropriate.

    \item \textbf{Essential information:} Names, contact details,
    employers, degrees, dates, locations, and important qualifications
    remain visible after text extraction.

    \item \textbf{Technical-skill clarity:} Programming languages,
    frameworks, databases, cloud services, and software tools are written
    explicitly and use clear multi-word names.

    \item \textbf{Evidence and context:} Important skills are supported
    by truthful project, course, research, or work-experience
    descriptions rather than appearing only as an isolated keyword list.

    \item \textbf{Job relevance:} Terminology from the target job
    description is included only when it accurately describes the
    applicant's qualifications.

    \item \textbf{Manual verification:} Automated feedback is checked
    against the original resume, and reported missing or matched terms
    are reviewed by a person before revisions are made.
\end{enumerate}

The rubric reflects the distinction between visual appearance and the
underlying text structure observed in the layout experiment. A creative
resume is not automatically unsuitable, but it requires additional
checking when it uses columns, graphics, icons, text boxes, or
non-standard reading order. Students should preserve a visually polished
design only when its essential content remains accessible after
extraction.

\subsection{Computer Science Resume Optimization Guide}

The Computer Science Resume Optimization Guide converts the
technical-skill and job-description findings into a repeatable review
process:

\begin{enumerate}
    \item Select a specific job description and identify the programming
    languages, tools, frameworks, databases, cloud platforms, and
    technical methods that it explicitly requires.

    \item Compare those requirements with the applicant's genuine
    coursework, projects, research, employment, and independent
    experience. Terms should be added only when they are accurate.

    \item Preserve complete multi-word technology names, such as
    ``Amazon Web Services'' or ``Microsoft SQL Server.'' When useful,
    include both a widely recognized abbreviation and its expanded form,
    such as ``Amazon Web Services (AWS).''

    \item Place core competencies in a clearly labeled Technical Skills
    section while also demonstrating important skills in project or
    experience descriptions. Context helps distinguish a genuine
    technical competency from an unrelated word match.

    \item Use consistent spelling and recognizable terminology, but do
    not repeat keywords unnaturally. Keyword stuffing may reduce
    readability and does not demonstrate proficiency.

    \item Export the completed resume as a text-readable PDF. Copy its
    contents into a plain text editor and inspect the reading order,
    missing characters, section boundaries, and multi-word skills.

    \item Treat automated match reports as diagnostic suggestions.
    Compare their output with the source resume and job description
    instead of relying on one extraction method or proprietary score.
\end{enumerate}

The guide is supported by the finding that no evaluated extraction method
recovered every relevant term. GPT performed slightly better for the 18
leading industry keywords, while the deterministic baseline recovered a
larger share of the complete 496-keyword vocabulary. Jobscan also
produced scores ranging from 21 to 50 across only four valid scans.
Consequently, the recommended process combines job-specific review,
machine-readability testing, and human judgment rather than optimizing
for one tool.

Together, the rubric and guide provide students with actionable checks
while preserving the study's cautious interpretation. Their purpose is
to reduce preventable extraction problems and improve the clarity of
truthful qualifications, not to manipulate hiring systems or promise a
particular employment result.

\section{Conclusion}

This research examined whether creative or complex resume layouts reduce
the accuracy of an ATS-style extraction and named-entity-recognition
pipeline and whether GPT-assisted and hybrid methods can improve
technical-skill recovery. The study combined a 20-resume PDF-layout
experiment, a manually reviewed 50-resume system comparison, full-corpus
processing of 1,020 computer-science-related resumes, an industry-keyword
analysis, and an exploratory Jobscan pilot.

The original directional hypothesis predicted that creative or complex
resumes would achieve lower strict F1 than standard resumes. The observed
results did not support that prediction. Under the original baseline,
creative resumes achieved an F1 of 0.329, compared with 0.223 for
standard resumes. After GPT-assisted reconstruction, the corresponding
scores increased to 0.530 and 0.375. Overall entity-level F1 increased
from 0.270 under the baseline to 0.447 under GPT reconstruction.
However, the resumes were not controlled versions of identical content,
and the sample contained only 10 documents in each layout group.
Consequently, these results do not establish that creative layouts are
generally easier or that layout has no effect. They show that visual
category alone did not explain performance in this sample.

The larger technical-skill evaluation demonstrated the value of
contextual extraction while also revealing the tradeoffs of hybrid
systems. On the manually reviewed 50-resume sample, Generalized V4
achieved a strict F1 of 0.276, GPT-only achieved 0.773, the raw hybrid
achieved 0.754, and the cross-validation-filtered hybrid achieved the
highest combined strict F1 of 0.782. The raw union increased recall but
introduced enough false positives to reduce its overall F1. Filtering
recovered much of that lost precision while retaining recall above the
GPT-only system.

The hybrid improvement was not uniform across dataset sources. GPT-only
remained slightly stronger for Jithin and Saugata, while the filtered
hybrid improved Snehaan strict F1 from 0.691 to 0.721. Bootstrap analysis
estimated a combined filtered-hybrid advantage of approximately 0.009,
but its 95\% confidence interval slightly crossed zero. The Snehaan
difference remained positive throughout its confidence interval. These
findings demonstrate that source-specific evaluation is necessary and
that an aggregate improvement should not be treated as evidence of
universal system superiority.

Across the complete 1,020-resume corpus, GPT produced 20,756 normalized
skill predictions. The raw V4--GPT union contained 25,238 predictions,
and the consolidated filter removed 2,502 V4-only predictions, producing
22,736 filtered-hybrid predictions. Because the complete corpus did not
have Gold Standard annotations, these totals measure prediction volume
and scalability rather than accuracy. Similarly, the industry-keyword
and Jobscan analyses provided external context but did not identify one
method or commercial score as universally reliable.

The practical conclusion is that students should focus on
machine-readable structure rather than relying on absolute claims about
resume appearance. Important content should remain selectable, text
extraction should preserve a logical reading order, technical skills
should be written explicitly and truthfully, and automated feedback
should be manually reviewed. A Resume Self-Audit Rubric and Computer
Science Resume Optimization Guide were developed to convert these
findings into cautious, actionable guidance.

Future research should test identical resume content across controlled
templates, use multiple independent annotators, validate the hybrid filter
on a newly annotated external dataset, and separate PDF text-extraction
errors from semantic entity-recognition errors. Additional work could
also evaluate other document parsers and commercial systems under
privacy-aware and ethically reviewed conditions. Overall, the study shows
that contextual models can substantially improve technical-skill
recovery, but reliable resume analysis still requires validated
evaluation, dataset-specific reporting, transparent system design, and
human oversight.

\section*{Acknowledgment}

The author thanks Dr. Sushil Poudel for his guidance, feedback, and
mentorship throughout this research. The author also acknowledges the
support of Caldwell University's STEM Advance program and the Department
of Computer Science, which provided the resources and research
environment necessary to complete this project.

\bibliographystyle{IEEEtran}
\bibliography{references}

\end{document}
